\documentclass[12pt]{article}
\usepackage[utf8]{inputenc}
\usepackage[T1]{fontenc}
\usepackage{hyperref}
\title{Yue Chinese Language Resources for K-12 Curriculum Development}
\author{Compiled by Claude Sonnet 5.5 (claude-sonnet-5-5) for Andr\'as Kornai}
\date{}
\begin{document}
\maketitle

\section{Introduction}
Yue Chinese (Cantonese; yue) has tens of millions of speakers in Guangdong, Hong Kong and Macao. The searches found mostly dictionaries, lexicons and materials for adult learners. They did not find a school-level curriculum for Hong Kong or Guangdong.

\section{Primary Educational Resources}
Course-style textbooks: Wikibooks Cantonese, a 30-lesson course with about 950 words of vocabulary (\url{https://en.wikibooks.org/wiki/Cantonese/How_To_Use_This_Textbook}); Structured Cantonese Courses (\url{https://cantonese.hk/courses}); UC Berkeley course list, which names the Commercial Press textbooks (\url{https://chinese.berkeley.edu/node/3760}).

\section{Linguistic and Reference Materials}
words.hk dataset, with over 53,000 entries (\url{https://repository.eduhk.hk/en/publications/wordshk-a-comprehensive-cantonese-dictionary-dataset-with-definit-2/}, \url{https://aclanthology.org/2022.dclrl-1.7}); Rime-Cantonese Jyutping lexicon (\url{https://catalog.ldc.upenn.edu/LDC2022L01}); Jyut Dictionary (\url{https://jyutdictionary.com/}); Yue Phonetic Database (\url{https://cjk.org/data/chinese/nlp/yue-phonetic-database}). Scholarship: Routledge, \emph{The Learning and Teaching of Cantonese as a Second Language} (\url{https://routledge.com/The-Learning-and-Teaching-of-Cantonese-as-a-Second-Language/Lee/p/book/9780367516321}); a review of Hong Kong Cantonese textbooks (\url{https://papers.iafor.org/submission36128/}).

\section{Religious and Community Resources}
None were found in the search results.

\section{Digital and Audio Resources}
Jyut Dictionary and words.hk, above.

\section{Curriculum Development Context}
Jyutping, from the Linguistic Society of Hong Kong, is the standard romanization (\url{https://simple.wikipedia.com/wiki/Jyutping}). The Education Bureau's Chinese Language curriculum was not found and needs a separate search.

\end{document}